% !TeX root = ../../Book3.tex
\chapter{椭圆函数}
\label{b3:ch:22}
\BookChapterNumber{b3:ch:22}
% 本章摘要（不计入编号）
\begin{chapterabstract}
双周期性将复平面的局部亚纯理论放到紧环面上。本章先用对边积分约束零点、极点和留数，再通过正规收敛格和构造 Weierstrass 函数。其微分方程给出光滑三次曲线，函数的次数与半周期分析进一步证明环面和该曲线双全纯，并解释弦切群律。
\end{chapterabstract}

% 本章目标（不计入编号）
\begin{learninggoals}
能够验证格和的收敛和周期性，计算椭圆函数的零极点总数及除数限制；掌握 \(\wp\) 的微分方程、分歧点和缩放规律；能从复参数说明三次曲线的群律与函数域。
\end{learninggoals}

一个整函数若同时沿两条独立方向重复，就会被困在一个紧平行四边形中的模界内。要得到新的函数，必须允许极点；但这些极点又受到留数和为零的限制。Weierstrass 的构造从一个无留数的二阶极点出发，将它复制到全格，给出这些限制的最小实例。所需基础为留数定理、辐角原理和正规级数；Chandrasekharan 的《Elliptic Functions》\cite{Chandrasekharan1985Elliptic}可用于继续比较格、椭圆函数与模函数三种表述。

% !TeX root = ../../Book3.tex
\section{格与双周期函数}
\label{b3:sec:2201}

指数函数把一个平移周期变成单值函数，两个实线性无关的周期则把平面压成紧曲面。紧性迫使非常值函数具有极点，因此双周期理论自然属于亚纯函数，而非整函数。

\subsection{\texorpdfstring{\(\mathbb C\)}{C} 中的格}
\label{b3:subsec:220101}

\begin{definition}\label{b3:def:22-lattice}
若 \(\omega_1,\omega_2\in\mathbb C\) 在实数域上线性无关，则离散加法子群
\(\Lambda=\mathbb Z\omega_1+\mathbb Z\omega_2\) 称为复平面中的\term[ge]{格}[lattice]。有序生成元称为格的一组基；通常调整方向使 \(\operatorname{Im}(\omega_2/\omega_1)>0\)。
\end{definition}
实线性映射 \((x,y)\mapsto x\omega_1+y\omega_2\) 可逆，故存在 \(c,C>0\)，使
\(c\sqrt{m^2+n^2}\le|m\omega_1+n\omega_2|\le C\sqrt{m^2+n^2}\)。因此格点没有有限聚点，半径 \(R\) 内有 \(O(1+R^2)\) 个格点，且
\[
\sum_{0\ne\omega\in\Lambda}|\omega|^{-p}<\infty\quad(p>2).
\]
最后结论可按整数方框壳 \(\max(|m|,|n|)=j\) 求和；该壳有 \(8j\) 个点，故归结于 \(\sum j^{1-p}\)。这条估计将负责椭圆级数的正规收敛。

\subsection{基本平行四边形}
\label{b3:subsec:220102}

集合 \(P=\{s\omega_1+t\omega_2:0\le s,t<1\}\) 是\term[jiben-pingxing-sibianxing]{基本平行四边形}[fundamental parallelogram]。每个复数模 \(\Lambda\) 有唯一一个代表落在 \(P\)；闭平行四边形上对边点须按平移等同。

\ProseParagraphBreak

积分时常将整个平行四边形平移一个小量，使边界不经过给定亚纯函数的零极点。这样的平移确实存在：在固定有界范围内只有有限多个零极点，需要避开的平移参数是一族有限线段，其并集不含开集。计算最终只取决于模格的点，不取决于平行四边形的起点。

\subsection{双周期函数}
\label{b3:subsec:220103}

若亚纯函数 \(f\) 满足 \(f(z+\omega)=f(z)\) 对全部 \(\omega\in\Lambda\) 成立，就称它关于 \(\Lambda\) 双周期。检验两组生成元足够，因正向平移、反向平移和有限次复合给出全部格平移。

\ProseParagraphBreak

周期群可能比指定的格更大。例如相对于 \(\Lambda\) 的函数也可相对于子格 \(2\Lambda\) 讨论，但后一基本区域中极点数变成原来的四倍。凡说“次数”或“模格的零点”，必须先固定采用哪个格；不能无说明地用函数的全部周期替换给定格。

\subsection{椭圆函数}
\label{b3:subsec:220104}

上述亚纯双周期函数称为\term[tuoyuan-hanshu]{椭圆函数}[elliptic function]。常函数也包含在这个定义中。
\begin{proposition}\label{b3:prop:22-entire-constant}
没有非常值的整椭圆函数。
\end{proposition}
\begin{proof}
整函数在闭基本平行四边形上有界；周期性把这个界传到全平面。Liouville 定理使它为常数。
\end{proof}
仅有一个实周期时结论失败：\(e^{2\pi\mathrm{i}z}\) 是周期为一的非常值整函数，其商空间仍有无穷远端。两个独立周期使商空间紧致，才造成上述刚性。

\ProseParagraphBreak

离散性也能从周期本身反向检验。设非常值亚纯函数有非零周期 \(\omega_j\to0\)。在任一没有极点的小圆盘内，差商
\(\bigl(f(z+\omega_j)-f(z)\bigr)/\omega_j\) 恒为零，极限给 \(f'(z)=0\)。因此函数在一个开集上常值，亚纯恒等定理使它处处常值，矛盾。所以非常值亚纯函数的周期群一定离散。特别地，三个实线性相关而产生稠密整数加法群的周期，不会产生比椭圆函数更自由的非常值函数。
\ProseParagraphBreak
所谓两个周期“独立”，必须是在实数域上独立。若两个非零实周期的比为无理数，它们生成的群在实轴稠密，上一论证同样迫使常数；若比为有理数，它们只生成一维离散周期群。只有二维实独立的情形给紧的基本区域。在判定题目中的函数是否真正双周期时，应先检查这个线性代数条件，而不是数出公式中出现了两个平移数值。
\ProseParagraphBreak
格和的收敛指数二也有几何含义。一个大圆中格点数不仅是 \(O(R^2)\)，还与面积同阶：用同样大小的基本平行四边形铺平面，圆内完全包含的格胞覆盖一个稍小圆，与圆相交的格胞则落在一个稍大圆。格胞面积固定，故计数介于 \(cR^2\) 与 \(CR^2\) 之间。于是未作相消的 \(\sum|\omega|^{-2}\) 在二维格上发散，而 \(\sum|\omega|^{-3}\) 收敛。Weierstrass 级数中的减项正是跨过这个临界指数所必需的修正。

% !TeX root = ../../Book3.tex
\section{椭圆函数的基本结构}
\label{b3:sec:2202}

平行四边形的对边积分相消，把局部极点资料转成全局限制。以下各积分均选取无零极点的边界，并把零极点按模格的重数计数。

\subsection{零点与极点}
\label{b3:subsec:220201}

\begin{proposition}\label{b3:prop:22-zero-pole-count}
非恒零椭圆函数在一个基本区域内的零点总重数等于极点总阶数。
\end{proposition}
\begin{proof}
对数导数 \(f'/f\) 仍双周期，对边方向相反，所以边界积分为零。辐角原理将积分除以 \(2\pi\mathrm{i}\) 识别为零点数减极点数。商空间紧致保证两类点都有限；等式不包含恒零函数。
\end{proof}
对任意有限值 \(a\)，同一结论应用于 \(f-a\)，说明取值 \(a\) 的原像按重数计数等于极点总阶数。这个整数对 \(a\) 不变，给出椭圆函数作为覆盖映射的次数。

\subsection{留数和}
\label{b3:subsec:220202}

\begin{proposition}\label{b3:prop:22-residue-sum}
椭圆函数在一个基本区域内的留数之和为零。
\end{proposition}
\begin{proof}
将上一证明的被积函数改成 \(f\)，用留数定理而不是辐角原理；对边仍因周期性相消。
\end{proof}
因此一个基本区域中不能恰有一个一阶极点，因为一阶极点的留数非零。但一个二阶极点可以出现，其留数可以为零。下一节的 \(\wp\) 函数正是次数二的最简单实例。两个一阶极点的留数必须相反；这个条件将允许用两个平移的对数导数相减来构造它。

\subsection{除数关系}
\label{b3:subsec:220203}

\begin{proposition}\label{b3:prop:22-divisor-necessary}
设非恒零椭圆函数的零点代表为 \(a_j\)，极点代表为 \(b_k\)，均按重数重复。则
\(\sum_j a_j-\sum_k b_k\in\Lambda\)。
\end{proposition}
\begin{proof}
令 \(h=f'/f\)。留数定理给差值等于 \((2\pi\mathrm{i})^{-1}\int_{\partial P}zh(z)\,\mathrm{d}z\)。将对边配对，得到两个基周期乘以沿另两条边的 \(h\) 积分，并附相应正负号。每一条这样的边，其端点处 \(f\) 值相同；沿边连续选对数，终点与起点对数相差 \(2\pi\mathrm{i}\) 的整数倍。因此所得差是 \(\omega_1,\omega_2\) 的整数线性组合。换代表只改变格元素，不影响结论。
\end{proof}
这一条件比零极点数相等更强。例如指定单零点 \(a\) 和单极点 \(b\) 会要求 \(a-b\in\Lambda\)，这时零极点在商空间上相同并相互抵消，所以次数一不可能出现。给定除数的存在性则还需要构造函数，不能从必要条件倒推；下文先建立足以进行这种构造的 \(\wp\) 及其积分。

\subsection{椭圆函数的次数}
\label{b3:subsec:220204}

若非常值椭圆函数在基本区域有总阶数 \(d\) 的极点，就称 \(d\) 为它相对于 \(\Lambda\) 的\term[tuoyuan-hanshu-cishu]{次数}[degree of an elliptic function]。局部展开 \(f(z)-a=c(z-z_0)^m+\cdots\) 中的 \(m\) 是该纤维点的重数；极点处改用 \(1/f\) 计算。

\ProseParagraphBreak

在非分歧值附近，各个局部逆分支连续存在，故原像个数恰为 \(d\)。分歧值上不同逆分支可能合并，但总重数不变。这样，辐角原理的整数不只是记账工具，而是刻画了映射的每条纤维。次数二的偶函数预示了一对原像 \(z,-z\)，它们只在半周期点合并。

% !TeX root = ../../Book3.tex
\section{Weierstrass \texorpdfstring{\(\wp\)}{℘} 函数}
\label{b3:sec:2203}

直接相加 \((z-\omega)^{-2}\) 的常数尾项只有临界的二次衰减。减去 \(\omega^{-2}\) 后，衰减提高一次，得到可在紧集上求导的格和。

\subsection{定义与收敛性}
\label{b3:subsec:220301}

\begin{definition}\label{b3:def:22-wp}
对格 \(\Lambda\)，定义\term[weierstrass-wp-hanshu]{Weierstrass 椭圆函数}[Weierstrass elliptic function]
\[
\wp(z)=\frac1{z^2}+\sum_{0\ne\omega\in\Lambda}
\left(\frac1{(z-\omega)^2}-\frac1{\omega^2}\right).
\]
\end{definition}
\symidx{\(\wp(z;\Lambda)\)：Weierstrass 椭圆函数}
若 \(|z|\le R\)、\(|\omega|>2R\)，括号内的模不超过 \(C_R|\omega|^{-3}\)。因此避开格点的紧集上正规收敛。在每个格点恰有一个二阶极点，主部为到该格点距离的倒数平方。把 \(\omega\) 与 \(-\omega\) 配对证明 \(\wp\) 为偶函数，且原点 Laurent 展开为
\[
\wp(z)=z^{-2}+3G_4z^2+5G_6z^4+7G_8z^6+\cdots,
\quad G_{2k}=\sum_{\omega\ne0}\omega^{-2k}.
\]
此式先在最短非零周期以内的圆盘成立，由几何级数展开和绝对收敛换序得到。

\subsection{\texorpdfstring{\(\wp'\)}{\AnalysisPdfWeierstrassP'}}
\label{b3:subsec:220302}

逐项求导后，\(\wp'(z)=-2\sum_{\omega\in\Lambda}(z-\omega)^{-3}\)。这个级数绝对正规收敛，平移后可直接重排，因此 \(\wp'\) 双周期。对每个基周期 \(\omega_j\)，差 \(\wp(z+\omega_j)-\wp(z)\) 为常数；取 \(z=-\omega_j/2\)，两点均非格点，偶性使常数为零。两基周期生成全部格，故 \(\wp\) 也双周期。

\ProseParagraphBreak

这一论证解释为何不直接在原级数中重排两个各自发散的二次格和。正规收敛只对扣除后的整项成立；先求导得到绝对收敛的三次格和，再积分回常数，是避免错误重排的关键。

\subsection{微分方程}
\label{b3:subsec:220303}

定义 \(g_2=60G_4\)、\(g_3=140G_6\)。
\begin{theorem}[Weierstrass 微分方程]\label{b3:thm:22-wp-equation}
函数 \(\wp\) 满足 \((\wp')^2=4\wp^3-g_2\wp-g_3\)。
\end{theorem}
\begin{proof}
写 \(\wp=z^{-2}+az^2+bz^4+O(z^6)\)，其中 \(a=3G_4\)、\(b=5G_6\)。展开得
\((\wp')^2-4\wp^3=-20az^{-2}-28b+O(z^2)\)。因此
\((\wp')^2-4\wp^3+g_2\wp+g_3\) 在原点可去且值为零。周期性使它在每个格点可去，其他点本来全纯；整椭圆函数必常数，原点值又使常数为零。
\end{proof}
再次求导并先在 \(\wp'\ne0\) 处约去，再用恒等定理，得到 \(\wp''=6\wp^2-g_2/2\)。它与初始主部一起递推确定全部 Laurent 系数，也给出下一节代数曲线的切线斜率。

\subsection{零点与极点}
\label{b3:subsec:220304}

模 \(\Lambda\) 的三个非零半周期为 \(\omega_1/2,\omega_2/2,(\omega_1+\omega_2)/2\)。奇性与周期性使 \(\wp'\) 在这三点为零。由于 \(\wp'\) 的唯一极点阶数为三，零点计数说明三点均为单零点，没有其他零点。

\ProseParagraphBreak

记它们的 \(\wp\) 值为 \(e_1,e_2,e_3\)。这些值两两不同：若两点具有同一值，则 \(\wp-e_j\) 在各点至少二阶消失，零点总数至少四，与其次数二矛盾。因此
\[
4X^3-g_2X-g_3=4(X-e_1)(X-e_2)(X-e_3),
\quad e_1+e_2+e_3=0.
\]
一般值的两个原像恰为 \(z,-z\)。若 \(\wp(z)=\wp(w)\)，\(\wp-\wp(w)\) 的次数二已经被 \(w,-w\) 耗尽；在半周期时这两个点合为一个二重点。这是后面证明解析表示单射的依据。

\ProseParagraphBreak

把 \(-\wp\) 积分一次会产生\term[weierstrass-zeta-hanshu]{Weierstrass zeta 函数}[Weierstrass zeta function]。为避免与后面的 Riemann zeta 函数混淆，记
\[
 Z_\Lambda(z)=\frac1z+\sum_{\omega\in\Lambda\setminus\{0\}}
 \left(\frac1{z-\omega}+\frac1\omega+\frac z{\omega^2}\right).
\]
在任意避开格点的紧集上，括号内为 \(O_K(|\omega|^{-3})\)，所以级数正规收敛且可逐项微分，得到 \(Z_\Lambda'=-\wp\)。对称配对 \(\omega,-\omega\) 说明它是奇函数。由于导数周期，
\(Z_\Lambda(z+\omega_j)-Z_\Lambda(z)=\eta_j\) 为常数，称为相应准周期。这两个准周期不可能同时为零，否则会得到只有一个单极点的椭圆函数，违反留数和。
\ProseParagraphBreak
对不经过格点、内部恰含一个格点的基本平行四边形积分 \(Z_\Lambda\)，留数定理给 \(2\pi\mathrm{i}\)。把上边与下边、右边与左边分别配对，准周期差给
\begin{equation}\label{b3:eq:22-legendre-relation}
 \eta_1\omega_2-\eta_2\omega_1=2\pi\mathrm{i}.
\end{equation}
这里基的方向为正。反转基顺序会同时反转围道方向和左侧符号。这一常数不是任意正规化，而是一个单位留数沿基本区域累计的结果。
\ProseParagraphBreak
再定义整函数
\[
 \Sigma_\Lambda(z)=z\prod_{\omega\ne0}
 \left(1-\frac z\omega\right)
 \exp\left(\frac z\omega+\frac{z^2}{2\omega^2}\right).
\]
每个尾因子的对数从三次项开始，所以格点三次幂和保证正规收敛。零点恰为格点且均单，原点导数为一，直接微分得 \(\Sigma_\Lambda'/\Sigma_\Lambda=Z_\Lambda\)。奇性也由 \(Z_\Lambda\) 的奇性确定：\(-\Sigma_\Lambda(-z)\) 与 \(\Sigma_\Lambda(z)\) 具有相同对数导数及原点导数，故相等。这就是\term[weierstrass-sigma-hanshu]{Weierstrass sigma 函数}[Weierstrass sigma function]，此处用大写符号避免与算术函数的约数和记号冲突。
\ProseParagraphBreak
对数导数的准周期性给商
\(\Sigma_\Lambda(z+\omega_j)/(e^{\eta_jz}\Sigma_\Lambda(z))\) 为非零常数。取 \(z=-\omega_j/2\)，用奇性及半周期不是格点，确定该常数为 \(-e^{\eta_j\omega_j/2}\)。于是
\[
 \Sigma_\Lambda(z+\omega_j)
 =-e^{\eta_j(z+\omega_j/2)}\Sigma_\Lambda(z).
\]
这条公式同时保留了指数和负号。若仅比较导数，会遗漏乘法常数；后面构造双周期商时，漏掉负号就可能产生错误的半周期。

\begin{theorem}[椭圆除数的实现]\label{b3:thm:22-divisor-realization}
设有限点列 \(a_1,\ldots,a_d\)、\(b_1,\ldots,b_d\) 按重数重复，且两列模格不相交。存在零点恰为前一列、极点恰为后一列的椭圆函数，当且仅当
\(\sum a_j-\sum b_j\in\Lambda\)。该函数相差一个非零常数唯一。
\end{theorem}
\begin{proof}
必要性已由基本区域的除数关系证明。对充分性，写 \(D=\sum a_j-\sum b_j=m\omega_1+n\omega_2\)，取 \(c=m\eta_1+n\eta_2\)，令
\[
 f(z)=e^{cz}\prod_{j=1}^d\frac{\Sigma_\Lambda(z-a_j)}{\Sigma_\Lambda(z-b_j)}.
\]
在平移 \(\omega_i\) 时，分子分母各有 \(d\) 个负号，恰好相消；总乘子为 \(e^{c\omega_i-\eta_iD}\)。由式~\eqref{b3:eq:22-legendre-relation}，指数在 \(i=1\) 时为 \(-2\pi\mathrm{i}n\)，在 \(i=2\) 时为 \(2\pi\mathrm{i}m\)，所以两乘子都为一。零极点由整乘积的单零点完全确定。若有两个实现函数，其商无零点、无极点且双周期，故为非零常数。
\end{proof}
构造中改变某个点的格代表，会改变指数因子与前面的常数，却不改变最终除数和按常数唯一的函数。给定零极点的总和条件因此是商环面上的条件，不能依赖选取平行四边形时的偶然代表。

% !TeX root = ../../Book3.tex
\section{椭圆曲线的解析表示}
\label{b3:sec:2204}

微分方程把一对亚纯函数放在一条代数曲线上。但“满足方程”尚不等于得到曲面的参数化，还需要验证无穷远点、局部坐标、单射和满射。

\subsection{三次方程}
\label{b3:subsec:220401}

把仿射方程 \(y^2=4x^3-g_2x-g_3\) 齐次化为
\[
E_\Lambda:\quad Y^2Z=4X^3-g_2XZ^2-g_3Z^3.
\]
在 \(Z=0\) 上只有一点 \(O=[0:1:0]\)。有限点的奇点须同时满足 \(y=0\)、\(12x^2-g_2=0\) 及曲线方程，等价于右侧三次多项式有重根。前节已经排除重根；在 \(O\) 处，对 \(Z\) 的偏导为一，也非奇点。因此这是一个光滑紧复曲线。

\subsection{判别式}
\label{b3:subsec:220402}

多项式 \(4X^3-g_2X-g_3\) 的无重根条件等价于
\(D=g_2^3-27g_3^2\ne0\)。这里把 \(D\) 称为该 Weierstrass 方程的\term[tuoyuan-quxian-panbieshi]{判别式}[discriminant]；有的代数规范另带固定常数，使用具体公式时应核对归一化。

\ProseParagraphBreak

换格为 \(a\Lambda\) 时，由绝对收敛格和可得
\[
\wp(az;a\Lambda)=a^{-2}\wp(z;\Lambda),\quad
g_2(a\Lambda)=a^{-4}g_2(\Lambda),\quad g_3(a\Lambda)=a^{-6}g_3(\Lambda).
\]
因此判别式的缩放次数为十二。这个整数将成为下一章模形式权十二的来源，并不是任意选择的指数。

\subsection{\texorpdfstring{\(\mathbb C/\Lambda\)}{C/\AnalysisPdfCapitalLambda}}
\label{b3:subsec:220403}

商空间 \(\mathbb C/\Lambda\) 由平行四边形粘合对边得到。每个小于半个最短周期的圆盘都与其非平凡格平移不相交，其投影给出复坐标。坐标变换为平移，因而全纯；闭平行四边形覆盖商空间，故它紧致。

\ProseParagraphBreak

对 \(z\notin\Lambda\) 定义 \(\Phi(z)=[\wp(z):\wp'(z):1]\)，对格点定义 \(\Phi(z)=O\)。在原点附近除以 \(\wp'\)，坐标为
\([\wp/\wp':1:1/\wp']=[-z/2+O(z^5):1:-z^3/2+O(z^7)]\)，所以该赋值全纯，且在原点为浸入。它按周期下降为商曲面上的全纯映射。

\subsection{解析统一化}
\label{b3:subsec:220404}

\begin{theorem}\label{b3:thm:22-uniformization}
映射 \(\Phi:\mathbb C/\Lambda\to E_\Lambda\) 是双全纯映射。
\end{theorem}
\begin{proof}
若两个有限点具有同一 \(x\) 坐标，则对应参数为 \(z\) 与 \(-z\)。它们的 \(y\) 坐标相反；若还相等，则 \(\wp'(z)=0\)，于是 \(z\) 为半周期，\(z\equiv-z\pmod\Lambda\)。无穷远点只来自格点，故单射。对任意 \((x,y)\in E_\Lambda\)，次数二保证 \(\wp\) 取到 \(x\)；微分方程给导数值为两个平方根之一，换 \(z\) 为 \(-z\) 即得到指定 \(y\)，所以满射。
在 \(\wp'\ne0\) 处可用 \(x\) 作局部坐标；在半周期处，\(\wp'\) 的零点为单零点，故 \(y\) 是局部坐标；在 \(O\) 处用上面 \(\wp/\wp'\) 的非零一阶项。逆函数定理在三种情形都给全纯局部逆，故全局逆也全纯。
\end{proof}
这里证明的是由格构造的曲线的解析表示。反过来从任意光滑复三次曲线恢复周期格，需要积分全纯微分并核对周期；这一一般方向可结合\appref{b3:app:A}的 Riemann 曲面理论继续学习，不能仅由上述显式代入推出。

\subsection{群律的解释}
\label{b3:subsec:220405}

通过 \(\Phi\) 将商群加法搬到曲线上，\(O\) 为单位元，\((x,y)\) 的逆为 \((x,-y)\)。其几何意义可由除数关系读出：一般直线与三次曲线交于三点，按重数记为 \(P,Q,R\)，相应线性函数在环面上只有原点三阶极点，因此三个参数之和模格为零。故 \(P+Q\) 是第三交点关于 \(x\) 轴的对称点。

\ProseParagraphBreak

当两点的横坐标不同时，直线斜率为 \(m=(y_2-y_1)/(x_2-x_1)\)，代入三次方程并比较根的和得
\[
x(P+Q)=m^2/4-x_1-x_2.
\]
对应纵坐标取第三交点的负值。两点合并时斜率改为 \((12x_1^2-g_2)/(2y_1)\)；纵坐标为零的切线情形给 \(2P=O\)。这些公式来自已经存在的群结构，结合律由参数相加自动保证，无需再对繁长有理式作三重展开。

\ProseParagraphBreak

还可确定整个椭圆函数域。若 \(f\) 为偶椭圆函数，它在 \(\wp\) 的每条纤维上取同一值，因此下降为球面上的亚纯函数，必为 \(\wp\) 的有理函数。在半周期处，局部参数变换 \(u\mapsto-u\) 使偶 Laurent 展开只含 \(u^{2k}\)，保证下降函数确实亚纯。对一般 \(f\)，偶部 \((f(z)+f(-z))/2\) 为 \(R(\wp)\)，奇部除以 \(\wp'\) 后也为偶函数，故
\[
f(z)=R(\wp(z))+\wp'(z)S(\wp(z))
\]
其中 \(R,S\) 有理。这说明曲线方程已经包含全部椭圆函数，而不只是两个偶然选择的函数。

\ProseParagraphBreak

三次曲线上的微分也可以通过参数直接理解。在非分歧点，\(x=\wp(z)\)、\(y=\wp'(z)\)，所以
\[
 \frac{\mathrm{d}x}{y}=\mathrm{d}z.
\]
在半周期处，两边仍全纯：若 \(z=z_0+t\)，则 \(x=e_j+\tfrac12\wp''(z_0)t^2+O(t^3)\)、\(y=\wp''(z_0)t+O(t^2)\)，分子与分母的一阶因子相消，商趋于 \(\mathrm{d}t\)。在无穷远，\(x=z^{-2}+O(z^2)\)、\(y=-2z^{-3}+O(z)\)，又得 \(\mathrm{d}x/y=(1+O(z^4))\mathrm{d}z\)。因此这个看似含分母的微分在整个紧曲线上无极点、无零点。
\ProseParagraphBreak
环面上的任意全纯微分提升到平面后写成 \(h(z)\mathrm{d}z\)。对格平移不变要求 \(h\) 是整椭圆函数，故常数。因此全纯微分空间恰为一维，由 \(\mathrm{d}z\) 生成。沿两条基本闭曲线积分得到 \(\omega_1,\omega_2\)，周期格遂从微分中恢复出来。这个结论是在已知格环面上直接计算；一般紧 Riemann 曲面的微分空间维数需要附录 A 的进一步理论。
\ProseParagraphBreak
二重映射 \(\wp\) 的四个分歧点也由局部式看得清楚。三个非零半周期在有限值 \(e_1,e_2,e_3\) 上有局部参数平方；原点的二阶极点在目标无穷远坐标 \(1/\wp\) 下同样是平方。其他点导数非零，局部为双全纯。因而三次曲线写成 \(y^2=4(x-e_1)(x-e_2)(x-e_3)\) 时，选择一个 \(x\) 通常有两种符号的 \(y\)，绕一个单根则交换这两种选择。代数方程的平方根分支与环面上的半周期是同一现象的两种坐标描述。
\ProseParagraphBreak
矩形格能进一步给出实图像。设 \(\omega_1>0\)、\(\omega_2\) 纯虚且虚部正，格关于共轭不变，由级数有 \(\overline{\wp(z)}=\wp(\bar z)\)。实轴上 \(\wp\) 实值，从原点附近的正无穷开始下降；在 \((0,\omega_1/2)\) 上导数没有零点，所以直到半周期才达到第一个临界值。它必须是三个实根中的最大者，否则下降途中先遇到别的根就会先有导数零点。沿虚轴，\(\wp(\mathrm{i}t)\sim-t^{-2}\)，从负无穷上升，首次临界值同理是最小根。第三个半周期给中间根。三个值确为实数，因为共轭半周期与原半周期相差格点。由此不用数值画图，就能确定矩形格的三个分歧值顺序及实椭圆曲线的两支。
\ProseParagraphBreak
上述顺序依赖矩形格假设。对一般倾斜格，三个分歧值可以全为非实复数，谈“最大根”便没有意义；但无重根、四个二阶分歧点和全纯微分的结论仍成立。几何图像有助于选择积分路径，却不能替代复参数情况下的局部代数检验。

% 本章小结（不计入编号）
\begin{chaptersummary}
格的二维点数估计保证扣除后的级数收敛，对边相消则把局部 Laurent 数据变成全局整数关系。\(\wp\) 与 \(\wp'\) 的次数分别为二与三，它们的三次方程给出完整的环面解析表示。缩放格改变 \(g_2,g_3\) 的权，而保持相应曲线的复结构；下一章将研究这种变化。
\end{chaptersummary}


\BookExercises

% 八题自拟：补充格的对称性、主部构造及曲线计算。
\begin{exercise}\label{b3:ex:22-1}
证明格的另一组正向基恰由 \(\mathrm{SL}_2(\mathbb Z)\) 的整系数基变换给出。
\end{exercise}
\begin{solution}
两组基互为整系数线性组合，故变换矩阵及其逆均整，行列式为 \(\pm1\)。保持方向使行列式为正，即一。反之这样的矩阵有整逆，生成同一格并保持方向。
\end{solution}
\begin{exercise}\label{b3:ex:22-2}
设 \(f\) 只有模格原点处一个至多二阶极点。证明 \(f=A\wp+B\)。
\end{exercise}
\begin{solution}
留数和为零排除主部中的一阶项。取 \(A\) 为二阶项系数，则 \(f-A\wp\) 无极点，整椭圆函数定理使其为常数。
\end{solution}
\begin{exercise}\label{b3:ex:22-3}
对方格 \(\Lambda=\mathbb Z+\mathrm{i}\mathbb Z\) 证明 \(g_3=0\)，并求 \(\wp(\mathrm{i}z)\) 与 \(\wp(z)\) 的关系。
\end{exercise}
\begin{solution}
由于 \(\mathrm{i}\Lambda=\Lambda\)，缩放给 \(g_3=\mathrm{i}^{-6}g_3=-g_3\)，故零；同理 \(\wp(\mathrm{i}z)=-\wp(z)\)。判别式非零进一步给 \(g_2\ne0\)。
\end{solution}
\begin{exercise}\label{b3:ex:22-4}
对 \(\Lambda=\mathbb Z+e^{\pi\mathrm{i}/3}\mathbb Z\) 证明 \(g_2=0\)、\(g_3\ne0\)。
\end{exercise}
\begin{solution}
记 \(a=e^{\pi\mathrm{i}/3}\)，有 \(a^2=a-1\)，故 \(a\Lambda=\Lambda\)。于是 \(g_2=a^{-4}g_2\)，而 \(a^{-4}\ne1\)，所以 \(g_2=0\)。判别式非零排除 \(g_3=0\)。
\end{solution}
\begin{exercise}\label{b3:ex:22-5}
证明 \(\wp(z)-\wp(a)\) 的除数为 \((a)+(-a)-2(0)\)，并区分 \(2a\in\Lambda\) 的情形。
\end{exercise}
\begin{solution}
假定 \(a\notin\Lambda\)。偶性给两个零点。若它们不同，次数二使两者均单；若为半周期，两点合并，导数零且二阶导数非零，故是二重零点。唯一极点在原点二阶。
\end{solution}
\begin{exercise}\label{b3:ex:22-6}
设椭圆函数 \(f\) 在原点的主部为 \(az^{-4}+bz^{-3}+cz^{-2}\)，且模格没有其他极点。用 \(\wp,\wp',\wp''\) 表示它。
\end{exercise}
\begin{solution}
三个函数的主部分别为 \(z^{-2},-2z^{-3},6z^{-4}\)，均无一阶项。于是 \(f-a\wp''/6+b\wp'/2-c\wp\) 无极点，为常数。再用微分方程可将 \(\wp''\) 换为 \(6\wp^2-g_2/2\)。
\end{solution}
\begin{exercise}\label{b3:ex:22-7}
求环面上满足 \(nz=0\) 的点数，并说明 \(n=2\) 时它们在三次曲线上的位置。
\end{exercise}
\begin{solution}
点由 \((a\omega_1+b\omega_2)/n\) 表示，\(0\le a,b<n\)，实线性独立保证这些 \(n^2\) 个点互异。\(n=2\) 时为单位元与三个非零半周期，对应 \(O\) 和 \((e_j,0)\)。
\end{solution}
\begin{exercise}\label{b3:ex:22-8}
设 \(a\) 不是格点或半周期。证明
\[
 Z_\Lambda(z+a)-Z_\Lambda(z)-Z_\Lambda(a)
 =\frac12\frac{\wp'(z)-\wp'(a)}{\wp(z)-\wp(a)}.
\]
\end{exercise}
\begin{solution}
两边关于 \(z\) 都是椭圆函数。右侧可能的极点在 \(0,\pm a\)：在 \(a\) 处分子分母同阶消去，在 \(-a\) 处留数为一，在0处主部为 \(-1/z\)。左侧在 \(-a,0\) 的留数也为一、负一，故两边之差无极点而为常数。令 \(z\to0\)，左侧加 \(1/z\) 后趋零；右侧利用 \(\wp=z^{-2}+O(z^2)\)、\(\wp'=-2z^{-3}+O(z)\) 展开，加 \(1/z\) 后也趋零。常数因而为零。半周期情形可经参数极限另外解释，但题目假设避免同时处理合并极点。
\end{solution}
